\documentclass[12pt]{article}
\usepackage[utf8]{inputenc}
\usepackage[T1]{fontenc}
\usepackage{hyperref}
\title{Vunjo Language Resources for K-12 Curriculum Development}
\author{Compiled by Claude Sonnet 5 (claude-sonnet-5) for Andr\'as Kornai}
\date{}
\begin{document}
\maketitle
\section{Introduction}
Vunjo (Kivunjo; ISO 639-3 vun; Glottolog vunj1238) is a Bantu language of the Kilimanjaro region of Tanzania, treated as part of the Central Chaga group. Speaker estimates in the sources found vary widely (about 300,000 from a 1992 count to over 800,000 for the people group on Joshua Project).
\section{Primary Educational Resources}
No primers, textbooks or school materials were found in search results.
\section{Linguistic and Reference Materials}
No dictionary or grammar of Kivunjo appeared in the results. Only scholarly articles on its phonology and syntax were mentioned, without links. A collection of Chagga proverbs and sayings exists: \url{https://afriprov.tangaza.ac.ke/endangered-african-proverbs-collections-chagga-tanzania-sayings/}. A language summary is at \url{https://rosettapanglossia.longnow.org/index.php?title=Vun}.
\section{Religious and Community Resources}
Joshua Project lists the people group and scripture status: \url{https://joshuaproject.net/people_groups/15806/TZ}. Global Recordings Network lists a Kivunjo New Testament audio recording via Faith Comes By Hearing.
\section{Digital and Audio Resources}
Audio materials: \url{https://globalrecordings.net/en/language/vun}
\section{Curriculum Development Context}
No school use was found. Teaching materials would need to be built from the audio scripture, proverb collections and new documentation. Sources disagree on whether the Bible listing is in Kivunjo or Swahili, so check this before relying on it.
\end{document}